% GENERATED FILE. Edit canonical lessons, then run npm run generate:books.
% canonical-source-hash: fnv1a64:d3ce7a6d3d01ba19
% canonical-lessons: GU-C429-tijori, GU-C429-nana, GU-C429-calan, GU-C429-lutavu, GU-C429-parvadvu

\chapter{Safes and Currency}
\label{ch:gu-safes-and-currency}
\hlchaptermodality{429}

\begin{chapteropening}
\textbf{By the end of this chapter:} \emph{I can talk about keeping money safe at home and abroad.}

Complete the last lesson of chapter 429: I can talk about keeping money safe at home and abroad.
\end{chapteropening}

\section[\texorpdfstring{\gu{તિજોરી} (tijorī)}{તિજોરી (tijorī)}]{\gu{તિજોરી} (tijorī) — a safe; a strongbox}

\label{lesson:GU-C429-tijori}



\begin{quote}
Before the new one: say the Gujarati for a moneylender, then the Gujarati for to swindle.
\end{quote}

\subsection*{You'll want to know: \gu{તિજોરી}}
\textbf{\gu{તિજોરી}} — \emph{tijorī} — \textquotedblleft{}a safe; a strongbox\textquotedblright{}.

\begin{quote}
\gu{દુકાનદાર} \gu{રોકડ} \gu{તિજોરીમાં} \gu{રાખે} \gu{છે}.
\end{quote}

\begin{quote}
\emph{dukāndār rokaḍ tijorīmā̃ rākhe che.}
\end{quote}

\begin{quote}
The shopkeeper keeps the cash in a safe.
\end{quote}

A feminine noun.

\begin{cousinweb}
\textbf{\gu{તિજોરી}} (\emph{tijorī}) is usually traced to Portuguese \emph{tesouro}, \textquotedblleft{}a treasure, a treasury\textquotedblright{}: the place where the money is locked away.
\end{cousinweb}

\begin{grammarlens}[title={What you've built}]
One more word for getting around.
\end{grammarlens}

\subsection*{Your turn}
\begin{itemize}
\raggedright
  \item \emph{Say it:} \emph{tijorī}
  \item \emph{Say it:} \emph{tijorī}, once more
  \item \emph{Say it:} \emph{dukāndār rokaḍ tijorīmā̃ rākhe che.} — The shopkeeper keeps the cash in a safe.
  \item \emph{Recall:} say the Gujarati for a moneylender, then the Gujarati for to swindle, then say \emph{tijorī} again
\end{itemize}

\begin{tcolorbox}[breakable,colback=teal!4,colframe=teal!35!black,title={Before you move on}]
What does \gu{તિજોરી} mean? (\textquotedblleft{}A safe; a strongbox\textquotedblright{}.) Say it once more.
\end{tcolorbox}
\section[\texorpdfstring{\gu{નાણાં} (nāṇā̃)}{નાણાં (nāṇā̃)}]{\gu{નાણાં} (nāṇā̃) — money; funds}

\label{lesson:GU-C429-nana}



\begin{quote}
Before the new one: say the Gujarati for to swindle, then the Gujarati for a safe.
\end{quote}

\subsection*{You'll want to know: \gu{નાણાં}}
\textbf{\gu{નાણાં}} — \emph{nāṇā̃} — \textquotedblleft{}money; funds\textquotedblright{}.

\begin{quote}
\gu{દાદાએ} \gu{નવી} \gu{શાળા} \gu{માટે} \gu{નાણાં} \gu{આપ્યાં}.
\end{quote}

\begin{quote}
\emph{dādāe navī śāḷā māṭe nāṇā̃ āpyā̃.}
\end{quote}

\begin{quote}
Grandfather gave funds for a new school.
\end{quote}

A neuter noun, used in the plural.

\begin{cousinweb}
\textbf{\gu{નાણાં}} (\emph{nāṇā̃}) goes back to Sanskrit \emph{nāṇaka}, \textquotedblleft{}a coin\textquotedblright{}; CDIAL lists the Gujarati word, \textquotedblleft{}a coin, money\textquotedblright{} (CDIAL 7050). The word is plural.
\end{cousinweb}

\begin{grammarlens}[title={What you've built}]
One more word for getting around.
\end{grammarlens}

\subsection*{Your turn}
\begin{itemize}
\raggedright
  \item \emph{Say it:} \emph{nāṇā̃}
  \item \emph{Say it:} \emph{nāṇā̃}, once more
  \item \emph{Say it:} \emph{dādāe navī śāḷā māṭe nāṇā̃ āpyā̃.} — Grandfather gave funds for a new school.
  \item \emph{Recall:} say the Gujarati for to swindle, then the Gujarati for a safe, then say \emph{nāṇā̃} again
\end{itemize}

\begin{tcolorbox}[breakable,colback=teal!4,colframe=teal!35!black,title={Before you move on}]
What does \gu{નાણાં} mean? (\textquotedblleft{}Money; funds\textquotedblright{}.) Say it once more.
\end{tcolorbox}
\section[\texorpdfstring{\gu{ચલણ} (calaṇ)}{ચલણ (calaṇ)}]{\gu{ચલણ} (calaṇ) — a currency}

\label{lesson:GU-C429-calan}



\begin{quote}
Before the new one: say the Gujarati for a safe, then the Gujarati for funds.
\end{quote}

\subsection*{You'll want to know: \gu{ચલણ}}
\textbf{\gu{ચલણ}} — \emph{calaṇ} — \textquotedblleft{}a currency\textquotedblright{}.

\begin{quote}
\gu{પરદેશમાં} \gu{ત્યાંનું} \gu{ચલણ} \gu{વાપરવું} \gu{પડે} \gu{છે}.
\end{quote}

\begin{quote}
\emph{pardeśmā̃ tyā̃nũ calaṇ vāparvũ paḍe che.}
\end{quote}

\begin{quote}
Abroad you have to use the local currency.
\end{quote}

A neuter noun.

\begin{cousinweb}
\textbf{\gu{ચલણ}} (\emph{calaṇ}) is a tadbhava of Sanskrit \emph{calana}, \textquotedblleft{}moving, going\textquotedblright{}; CDIAL lists the Gujarati word, \textquotedblleft{}conduct; currency\textquotedblright{} (CDIAL 4716): money that goes round.
\end{cousinweb}

\begin{grammarlens}[title={What you've built}]
One more word for getting around.
\end{grammarlens}

\subsection*{Your turn}
\begin{itemize}
\raggedright
  \item \emph{Say it:} \emph{calaṇ}
  \item \emph{Say it:} \emph{calaṇ}, once more
  \item \emph{Say it:} \emph{pardeśmā̃ tyā̃nũ calaṇ vāparvũ paḍe che.} — Abroad you have to use the local currency.
  \item \emph{Recall:} say the Gujarati for a safe, then the Gujarati for funds, then say \emph{calaṇ} again
\end{itemize}

\begin{tcolorbox}[breakable,colback=teal!4,colframe=teal!35!black,title={Before you move on}]
What does \gu{ચલણ} mean? (\textquotedblleft{}A currency\textquotedblright{}.) Say it once more.
\end{tcolorbox}
\section[\texorpdfstring{\gu{લૂંટાવું} (lū̃ṭāvũ)}{લૂંટાવું (lū̃ṭāvũ)}]{\gu{લૂંટાવું} (lū̃ṭāvũ) — to be robbed}

\label{lesson:GU-C429-lutavu}



\begin{quote}
Before the new one: say the Gujarati for funds, then the Gujarati for a currency.
\end{quote}

\subsection*{You'll want to know: \gu{લૂંટાવું}}
\textbf{\gu{લૂંટાવું}} — \emph{lū̃ṭāvũ} — \textquotedblleft{}to be robbed\textquotedblright{}.

\begin{quote}
\gu{રાત્રે} \gu{રસ્તામાં} \gu{વેપારી} \gu{લૂંટાયો}.
\end{quote}

\begin{quote}
\emph{rātre rastāmā̃ vepārī lū̃ṭāyo.}
\end{quote}

\begin{quote}
The merchant was robbed on the road at night.
\end{quote}

Past (he …): \gu{તે} \gu{વેપારી} \gu{લૂંટાયો}, \emph{te vepārī lū̃ṭāyo}.

\begin{cousinweb}
\textbf{\gu{લૂંટાવું}} (\emph{lū̃ṭāvũ}) is the intransitive partner of \gu{લૂંટવું}, \textquotedblleft{}to rob\textquotedblright{}, which goes back to \emph{*luṭṭati}, \textquotedblleft{}plunders\textquotedblright{}, a form CDIAL reconstructs (CDIAL 11078): who robs, \gu{લૂંટવું}; who is robbed, \gu{લૂંટાવું}.
\end{cousinweb}

\begin{grammarlens}[title={What you've built}]
One more word for getting around.
\end{grammarlens}

\subsection*{Your turn}
\begin{itemize}
\raggedright
  \item \emph{Say it:} \emph{lū̃ṭāvũ}
  \item \emph{Say it:} \emph{lū̃ṭāvũ}, once more
  \item \emph{Say it:} \emph{rātre rastāmā̃ vepārī lū̃ṭāyo.} — The merchant was robbed on the road at night.
  \item \emph{Recall:} say the Gujarati for funds, then the Gujarati for a currency, then say \emph{lū̃ṭāvũ} again
\end{itemize}

\begin{tcolorbox}[breakable,colback=teal!4,colframe=teal!35!black,title={Before you move on}]
What does \gu{લૂંટાવું} mean? (\textquotedblleft{}To be robbed\textquotedblright{}.) Say it once more.
\end{tcolorbox}
\section[\texorpdfstring{\gu{પરવડવું} (parvaḍvũ)}{પરવડવું (parvaḍvũ)}]{\gu{પરવડવું} (parvaḍvũ) — to be affordable; to be within one's means}

\label{lesson:GU-C429-parvadvu}



\begin{quote}
Before the new one: say the Gujarati for a currency, then the Gujarati for to be robbed.
\end{quote}

\subsection*{You'll want to know: \gu{પરવડવું}}
\textbf{\gu{પરવડવું}} — \emph{parvaḍvũ} — \textquotedblleft{}to be affordable; to be within one's means\textquotedblright{}.

\begin{quote}
\gu{આટલી} \gu{મોંઘી} \gu{ગાડી} \gu{મને} \gu{નહીં} \gu{પરવડે}.
\end{quote}

\begin{quote}
\emph{āṭlī mõghī gāḍī mane nahī̃ paravḍe.}
\end{quote}

\begin{quote}
I will not be able to afford such an expensive car.
\end{quote}

Past (it …): \gu{તે} \gu{ભાડું} \gu{તેને} \gu{પરવડ્યું}, \emph{te bhāḍũ tene parvaḍyũ}.

\begin{cousinweb}
\textbf{\gu{પરવડવું}} (\emph{parvaḍvũ}) is a tadbhava of Sanskrit \emph{paripatati}, \textquotedblleft{}runs about\textquotedblright{}; CDIAL lists the Gujarati verb, \textquotedblleft{}to be suitable\textquotedblright{} (CDIAL 7841): what suits your purse.
\end{cousinweb}

\begin{grammarlens}[title={What you've built}]
One more word for getting around.
\end{grammarlens}

\subsection*{Your turn}
\begin{itemize}
\raggedright
  \item \emph{Say it:} \emph{parvaḍvũ}
  \item \emph{Say it:} \emph{parvaḍvũ}, once more
  \item \emph{Say it:} \emph{āṭlī mõghī gāḍī mane nahī̃ paravḍe.} — I will not be able to afford such an expensive car.
  \item \emph{Recall:} say the Gujarati for a currency, then the Gujarati for to be robbed, then say \emph{parvaḍvũ} again
\end{itemize}

\begin{tcolorbox}[breakable,colback=teal!4,colframe=teal!35!black,title={Before you move on}]
What does \gu{પરવડવું} mean? (\textquotedblleft{}To be affordable; to be within one's means\textquotedblright{}.) Say it once more.
\end{tcolorbox}
